% Options for packages loaded elsewhere
\PassOptionsToPackage{unicode}{hyperref}
\PassOptionsToPackage{hyphens}{url}
\PassOptionsToPackage{dvipsnames,svgnames,x11names}{xcolor}
%
\documentclass[
  10pt,
  a4paper,
]{article}
\usepackage{amsmath,amssymb}
\usepackage{iftex}
\ifPDFTeX
  \usepackage[T1]{fontenc}
  \usepackage[utf8]{inputenc}
  \usepackage{textcomp} % provide euro and other symbols
\else % if luatex or xetex
  \usepackage{unicode-math} % this also loads fontspec
  \defaultfontfeatures{Scale=MatchLowercase}
  \defaultfontfeatures[\rmfamily]{Ligatures=TeX,Scale=1}
\fi
\usepackage{lmodern}
\ifPDFTeX\else
  % xetex/luatex font selection
\fi
% Use upquote if available, for straight quotes in verbatim environments
\IfFileExists{upquote.sty}{\usepackage{upquote}}{}
\IfFileExists{microtype.sty}{% use microtype if available
  \usepackage[]{microtype}
  \UseMicrotypeSet[protrusion]{basicmath} % disable protrusion for tt fonts
}{}
\makeatletter
\@ifundefined{KOMAClassName}{% if non-KOMA class
  \IfFileExists{parskip.sty}{%
    \usepackage{parskip}
  }{% else
    \setlength{\parindent}{0pt}
    \setlength{\parskip}{6pt plus 2pt minus 1pt}}
}{% if KOMA class
  \KOMAoptions{parskip=half}}
\makeatother
\usepackage{xcolor}
\usepackage[margin=2.2cm]{geometry}
\usepackage{color}
\usepackage{fancyvrb}
\newcommand{\VerbBar}{|}
\newcommand{\VERB}{\Verb[commandchars=\\\{\}]}
\DefineVerbatimEnvironment{Highlighting}{Verbatim}{commandchars=\\\{\}}
% Add ',fontsize=\small' for more characters per line
\newenvironment{Shaded}{}{}
\newcommand{\AlertTok}[1]{\textcolor[rgb]{1.00,0.00,0.00}{\textbf{#1}}}
\newcommand{\AnnotationTok}[1]{\textcolor[rgb]{0.38,0.63,0.69}{\textbf{\textit{#1}}}}
\newcommand{\AttributeTok}[1]{\textcolor[rgb]{0.49,0.56,0.16}{#1}}
\newcommand{\BaseNTok}[1]{\textcolor[rgb]{0.25,0.63,0.44}{#1}}
\newcommand{\BuiltInTok}[1]{\textcolor[rgb]{0.00,0.50,0.00}{#1}}
\newcommand{\CharTok}[1]{\textcolor[rgb]{0.25,0.44,0.63}{#1}}
\newcommand{\CommentTok}[1]{\textcolor[rgb]{0.38,0.63,0.69}{\textit{#1}}}
\newcommand{\CommentVarTok}[1]{\textcolor[rgb]{0.38,0.63,0.69}{\textbf{\textit{#1}}}}
\newcommand{\ConstantTok}[1]{\textcolor[rgb]{0.53,0.00,0.00}{#1}}
\newcommand{\ControlFlowTok}[1]{\textcolor[rgb]{0.00,0.44,0.13}{\textbf{#1}}}
\newcommand{\DataTypeTok}[1]{\textcolor[rgb]{0.56,0.13,0.00}{#1}}
\newcommand{\DecValTok}[1]{\textcolor[rgb]{0.25,0.63,0.44}{#1}}
\newcommand{\DocumentationTok}[1]{\textcolor[rgb]{0.73,0.13,0.13}{\textit{#1}}}
\newcommand{\ErrorTok}[1]{\textcolor[rgb]{1.00,0.00,0.00}{\textbf{#1}}}
\newcommand{\ExtensionTok}[1]{#1}
\newcommand{\FloatTok}[1]{\textcolor[rgb]{0.25,0.63,0.44}{#1}}
\newcommand{\FunctionTok}[1]{\textcolor[rgb]{0.02,0.16,0.49}{#1}}
\newcommand{\ImportTok}[1]{\textcolor[rgb]{0.00,0.50,0.00}{\textbf{#1}}}
\newcommand{\InformationTok}[1]{\textcolor[rgb]{0.38,0.63,0.69}{\textbf{\textit{#1}}}}
\newcommand{\KeywordTok}[1]{\textcolor[rgb]{0.00,0.44,0.13}{\textbf{#1}}}
\newcommand{\NormalTok}[1]{#1}
\newcommand{\OperatorTok}[1]{\textcolor[rgb]{0.40,0.40,0.40}{#1}}
\newcommand{\OtherTok}[1]{\textcolor[rgb]{0.00,0.44,0.13}{#1}}
\newcommand{\PreprocessorTok}[1]{\textcolor[rgb]{0.74,0.48,0.00}{#1}}
\newcommand{\RegionMarkerTok}[1]{#1}
\newcommand{\SpecialCharTok}[1]{\textcolor[rgb]{0.25,0.44,0.63}{#1}}
\newcommand{\SpecialStringTok}[1]{\textcolor[rgb]{0.73,0.40,0.53}{#1}}
\newcommand{\StringTok}[1]{\textcolor[rgb]{0.25,0.44,0.63}{#1}}
\newcommand{\VariableTok}[1]{\textcolor[rgb]{0.10,0.09,0.49}{#1}}
\newcommand{\VerbatimStringTok}[1]{\textcolor[rgb]{0.25,0.44,0.63}{#1}}
\newcommand{\WarningTok}[1]{\textcolor[rgb]{0.38,0.63,0.69}{\textbf{\textit{#1}}}}
\setlength{\emergencystretch}{3em} % prevent overfull lines
\providecommand{\tightlist}{%
  \setlength{\itemsep}{0pt}\setlength{\parskip}{0pt}}
\setcounter{secnumdepth}{-\maxdimen} % remove section numbering
\usepackage{microtype}
\usepackage{longtable,booktabs,array}
\usepackage{xcolor}
\usepackage{fancyhdr}
\usepackage{fvextra}
\DefineVerbatimEnvironment{Highlighting}{Verbatim}{breaklines,breakanywhere,commandchars=\\\{\}}
\pagestyle{fancy}
\fancyhf{}
\fancyhead[L]{DDTA Research Methodology}
\fancyhead[R]{BA5 / Controlled Authoring}
\fancyfoot[C]{\thepage}
\setlength{\headheight}{14pt}
\emergencystretch=3em
\ifLuaTeX
  \usepackage{selnolig}  % disable illegal ligatures
\fi
\usepackage{bookmark}
\IfFileExists{xurl.sty}{\usepackage{xurl}}{} % add URL line breaks if available
\urlstyle{same}
\hypersetup{
  pdftitle={DDTA BA5-T3 canonical registry, controlled-authoring and no-synonym closure review - R1},
  colorlinks=true,
  linkcolor={blue},
  filecolor={Maroon},
  citecolor={Blue},
  urlcolor={blue},
  pdfcreator={LaTeX via pandoc}}

\title{DDTA BA5-T3 canonical registry, controlled-authoring and
no-synonym closure review - R1}
\author{}
\date{}

\begin{document}
\maketitle

\section{DDTA BA5-T3 canonical registry, controlled-authoring and
no-synonym closure review -
R1}\label{ddta-ba5-t3-canonical-registry-controlled-authoring-and-no-synonym-closure-review---r1}

\textbf{Status:} CLOSED / PASS

\textbf{Repository baseline reviewed:}
\texttt{8d8dae5f7c28d83b70cbdea090028e4ec0f93571}

\textbf{Microstep:}
\texttt{BA5-T3\ -\ canonical\ registry,\ controlled-authoring\ and\ no-synonym\ closure\ review}

\textbf{Scope boundary:} this review closes BA5 only if the integrated
T1/T2 contract survives adversarial pressure. It does not execute BA6,
does not add a method schema and does not introduce ThreatForge
implementation semantics.

\subsection{1. Question under test}\label{1-question-under-test}

The closure question is:

\begin{quote}
Can DDTA, for its semantically operative authoring surface, rely on
exact canonical referent names and exact domain-scoped semantic tokens,
with governed extension and tool enforcement, without normative
synonym/alias machinery and without reopening BA0-BA4?
\end{quote}

The burden of proof is on additional lexical machinery. Convenience is
not sufficient evidence.

\subsection{2. Inputs integrated by T3}\label{2-inputs-integrated-by-t3}

T3 integrates:

\begin{itemize}
\tightlist
\item
  BA5-T1 exact canonical referent naming;
\item
  BA5-T2 domain-scoped operator/role/kind/value registries;
\item
  BA2 closed operator and operator-scoped role semantics;
\item
  BA3 cross-baseline identity continuity;
\item
  BA4 shared-rendering and method-owned interpretation separation; and
\item
  the portable-by-construction documentation assumption already fixed by
  the project.
\end{itemize}

The review attacks the combined contract rather than repeating the
individual T1/T2 trials.

\subsection{3. Falsification matrix}\label{3-falsification-matrix}

The integrated attacks and dispositions are:

\begin{itemize}
\tightlist
\item
  \textbf{C1 - referent alias:} second normative identifier required
  -\textgreater{} \textbf{PASS; alias rejected}.
\item
  \textbf{C2 - operator alias:} \texttt{send} required beside
  \texttt{transfer} -\textgreater{} \textbf{PASS}.
\item
  \textbf{C3 - role alias:} \texttt{src} required beside \texttt{source}
  -\textgreater{} \textbf{PASS}.
\item
  \textbf{C4 - semantic-kind alias:} \texttt{repository} required beside
  \texttt{store} for the same meaning -\textgreater{} \textbf{PASS}.
\item
  \textbf{C5 - controlled-value alias:} \texttt{owner} required beside
  \texttt{ownership} -\textgreater{} \textbf{PASS}.
\item
  \textbf{C6 - same token in two domains:} flat collision makes lookup
  unusable -\textgreater{} \textbf{PASS WITH DOMAIN-SCOPE}.
\item
  \textbf{C7 - governed rename B0 to B1:} name must become BA
  identity/alias history -\textgreater{} \textbf{PASS WITH BA3 RETAIN}.
\item
  \textbf{C8 - registry evolution:} historical semantics cannot be
  replayed -\textgreater{} \textbf{PASS WITH IMMUTABLE REVISION}.
\item
  \textbf{C9 - ordinary human prose:} all prose must become controlled
  language -\textgreater{} \textbf{PASS WITH OPERABILITY BOUNDARY}.
\item
  \textbf{C10 - incompatible projections:} local taxonomy must leak into
  BA -\textgreater{} \textbf{PASS}.
\item
  \textbf{C11 - method-owned aggregation:} shared referent must be
  renamed -\textgreater{} \textbf{PASS WITH BA4 BOUNDARY}.
\item
  \textbf{C12 - typed literal:} every local value must become registry
  vocabulary -\textgreater{} \textbf{PASS; literal remains local}.
\item
  \textbf{C13 - genuine new semantic kind:} no extension path exists
  without synonyms -\textgreater{} \textbf{PASS WITH GOVERNED
  EXTENSION}.
\item
  \textbf{C14 - material new operator/role:} BA5 must silently change
  BA2 -\textgreater{} \textbf{PASS; routed to BA2 reopen}.
\item
  \textbf{C15 - tool/LLM alias suggestion:} suggestion must become truth
  -\textgreater{} \textbf{PASS; candidate only}.
\item
  \textbf{C16 - translated/descriptive presentation:} second semantic
  identifier is unavoidable -\textgreater{} \textbf{PASS; binding
  remains canonical}.
\end{itemize}

No trial produces a material counterexample that forces normative
synonym support.

\subsection{4. C1 - referent alias
attack}\label{4-c1---referent-alias-attack}

Start with:

\begin{Shaded}
\begin{Highlighting}[]
\NormalTok{canonical project referent:}
\NormalTok{  cameraIngresso}
\end{Highlighting}
\end{Shaded}

Attack:

\begin{Shaded}
\begin{Highlighting}[]
\NormalTok{D1 {-}\textgreater{} cameraIngresso}
\NormalTok{D2 {-}\textgreater{} entryCamera}
\NormalTok{human view {-}\textgreater{} Camera Ingresso}
\end{Highlighting}
\end{Shaded}

If all three intend the same shared referent, the additional identifiers
are not admitted as normative semantic names. The author/view uses
\texttt{cameraIngresso} for the semantic reference.

Descriptive prose can still say "entrance camera" around that reference.

\textbf{Result:} PASS. No alias registry is required.

\subsection{5. C2/C3 - operator and role alias
attacks}\label{5-c2c3---operator-and-role-alias-attacks}

Canonical binding:

\begin{Shaded}
\begin{Highlighting}[]
\NormalTok{transfer}
\NormalTok{  source      {-}\textgreater{} cameraIngresso}
\NormalTok{  destination {-}\textgreater{} recognitionProcessor}
\NormalTok{  content     {-}\textgreater{} recognitionCapture}
\end{Highlighting}
\end{Shaded}

Attacks:

\begin{Shaded}
\begin{Highlighting}[]
\NormalTok{send}
\NormalTok{src}
\NormalTok{dst}
\NormalTok{payload}
\end{Highlighting}
\end{Shaded}

None is needed as a normative key. A human authoring interface may
search for or suggest \texttt{transfer}, \texttt{source},
\texttt{destination}, \texttt{content}, but the persisted/accepted
semantic binding remains exact.

\textbf{Result:} PASS.

\subsection{6. C4/C5 - semantic-kind and controlled-value
aliases}\label{6-c4c5---semantic-kind-and-controlled-value-aliases}

Attacks:

\begin{Shaded}
\begin{Highlighting}[]
\NormalTok{repository  {-}\textgreater{} intended same meaning as store}
\NormalTok{owner       {-}\textgreater{} intended same responsibility meaning as ownership}
\end{Highlighting}
\end{Shaded}

When the intended meaning is already represented, the alternative token
is rejected as a second normative key.

If \texttt{repository} actually denotes a materially different
methodology-neutral concept, it is not treated as a synonym. It becomes
an evidence-backed candidate under the semantic-kind extension rule.

\textbf{Result:} PASS. Lexical difference does not decide semantic
difference.

\subsection{7. C6 - cross-domain token
collision}\label{7-c6---cross-domain-token-collision}

Attack:

\begin{Shaded}
\begin{Highlighting}[]
\NormalTok{referent canonicalName = source}
\NormalTok{transfer roleKey        = source}
\end{Highlighting}
\end{Shaded}

A flat global dictionary fails here, but T2 already refined the
hypothesis to domain-scoped resolution.

Lookup remains unambiguous because the two tokens resolve through
different domains/scopes.

\textbf{Result:} PASS WITH DOMAIN-SCOPE REFINEMENT retained. No semantic
collapse.

\subsection{8. C7 - cross-baseline
rename}\label{8-c7---cross-baseline-rename}

Attack:

\begin{Shaded}
\begin{Highlighting}[]
\NormalTok{B0  cameraIngresso}
\NormalTok{B1  cameraNord}
\end{Highlighting}
\end{Shaded}

If the independently reusable project meaning survives, BA3 can RETAIN
the same \texttt{BAReferent} identity while the current baseline changes
its canonical name.

Historical material does not need an alias bridge:

\begin{Shaded}
\begin{Highlighting}[]
\NormalTok{B0 artifacts {-}\textgreater{} cameraIngresso}
\NormalTok{B1 artifacts {-}\textgreater{} cameraNord}
\NormalTok{continuity   {-}\textgreater{} BA3 RETAIN}
\end{Highlighting}
\end{Shaded}

\textbf{Result:} PASS. Name remains baseline-scoped lexical metadata
rather than semantic identity.

\subsection{9. C8 - semantic registry
evolution}\label{9-c8---semantic-registry-evolution}

Attack: keep token \texttt{store} while silently changing its normative
definition.

This would destroy historical reproducibility.

The closed answer is immutable registry revision resolution:

\begin{Shaded}
\begin{Highlighting}[]
\NormalTok{SemanticKindRegistry@R1 / store}
\NormalTok{SemanticKindRegistry@R2 / store}
\end{Highlighting}
\end{Shaded}

If R2 changes meaning materially, the change is inspectable and
historical BA continues to resolve R1. A semantic change cannot be
hidden behind the same unversioned token.

\textbf{Result:} PASS WITH IMMUTABLE REVISION requirement retained.

\subsection{10. C9 - natural-language
pressure}\label{10-c9---natural-language-pressure}

A too-strong interpretation would require every sentence to use only
canonical verbs and nouns. That is unnecessary and would convert BA5
into a controlled-natural-language project.

The smallest sufficient boundary is:

\begin{Shaded}
\begin{Highlighting}[]
\NormalTok{human explanation}
\NormalTok{    may use natural language}

\NormalTok{machine{-}significant / governed semantic binding}
\NormalTok{    must use canonical registry key/name}
\end{Highlighting}
\end{Shaded}

For example "sends" in prose does not create a \texttt{send} operator
when the semantic binding is \texttt{transfer}.

\textbf{Result:} PASS WITH OPERABILITY BOUNDARY. This is a minimization,
not a relaxation of semantic consistency.

\subsection{11. C10/C11 - projection incompatibility
pressure}\label{11-c10c11---projection-incompatibility-pressure}

Two incompatible consumers may classify/render the same shared referent
differently:

\begin{Shaded}
\begin{Highlighting}[]
\NormalTok{FlowParticipant(cameraIngresso)}
\NormalTok{AssuranceSubject(cameraIngresso)}
\end{Highlighting}
\end{Shaded}

The projection-owned kinds differ while the shared referent name is
stable.

A genuinely method-owned aggregation may instead use:

\begin{Shaded}
\begin{Highlighting}[]
\NormalTok{external{-}ingress{-}exposure}
\end{Highlighting}
\end{Shaded}

because that label denotes a downstream method item rather than the
shared referent. BA4 trace and interpretation rules preserve the
boundary.

\textbf{Result:} PASS. Method terminology does not need to enter the BA
registry.

\subsection{12. C12 - literal inflation
attack}\label{12-c12---literal-inflation-attack}

Attack:

\begin{Shaded}
\begin{Highlighting}[]
\NormalTok{timeout = 30s}
\end{Highlighting}
\end{Shaded}

Registering every local value would turn the registry into a universal
value store. The controlled-value registry applies only to reusable
governed semantic distinctions.

\texttt{30s} remains a typed/local value unless independent reusable
semantic identity forces a \texttt{BAReferent} under BA1.

\textbf{Result:} PASS. Registry inflation is not required.

\subsection{13. C13/C14 - extension routing
pressure}\label{13-c13c14---extension-routing-pressure}

\subsubsection{New method-neutral semantic
kind}\label{new-method-neutral-semantic-kind}

A genuinely new kind may be admitted only after evidence, method-neutral
definition and review, producing a new immutable registry revision.

\subsubsection{New operator or role
semantic}\label{new-operator-or-role-semantic}

A materially new operator or changed role contract is not a lexical
addition. It reopens the smallest BA2 responsibility.

Therefore BA5 cannot become a backdoor for semantic expansion.

\textbf{Result:} PASS.

\subsection{14. C15 - tool/LLM authority
attack}\label{14-c15---toolllm-authority-attack}

Attack:

\begin{Shaded}
\begin{Highlighting}[]
\NormalTok{user writes: send}
\NormalTok{model confidence: 0.97 that send == transfer}
\end{Highlighting}
\end{Shaded}

Allowed behavior:

\begin{Shaded}
\begin{Highlighting}[]
\NormalTok{"send is not a registered operator; did you mean transfer?"}
\end{Highlighting}
\end{Shaded}

Rejected behavior:

\begin{Shaded}
\begin{Highlighting}[]
\NormalTok{silently persist transfer}
\NormalTok{or register send as an accepted synonym}
\end{Highlighting}
\end{Shaded}

The author/governance path decides accepted semantic binding. The model
remains assistance, not authority.

\textbf{Result:} PASS.

\subsection{15. C16 - translation/accessibility
pressure}\label{15-c16---translationaccessibility-pressure}

A human view may need explanatory or translated presentation.

This does not force a second normative semantic identifier. For example:

\begin{Shaded}
\begin{Highlighting}[]
\NormalTok{Entrance camera (\textasciigrave{}cameraIngresso\textasciigrave{})}
\end{Highlighting}
\end{Shaded}

may be presentation text while the underlying/shared semantic reference
remains \texttt{cameraIngresso}.

If a downstream method owns a translated/local item rather than the
shared referent, BA4 method-owned labeling applies.

\textbf{Result:} PASS. Translation support and semantic aliasing remain
distinct responsibilities.

\subsection{16. Cross-corpus closure
reasoning}\label{16-cross-corpus-closure-reasoning}

The two already reviewed project pressures remain compatible with the
closed BA5 boundary:

\begin{itemize}
\tightlist
\item
  facial-access cases use stable named referents plus the closed BA2
  operators/roles for transfer, correlation, service consumption,
  realization, responsibility and constraints;
\item
  order, WMS and provider cases use stable referent naming plus
  dependency, service, responsibility and classification semantics and
  governed project-specific mappings without requiring a second
  normative lexical key.
\end{itemize}

T3 does not claim arbitrary prose migration. It asks whether
intentionally analysis-ready governed documentation can preserve the
shared semantics using canonical bindings. The reviewed cases do not
falsify that claim.

A structurally different holdout remains reserved for BA6 integrated
regression, as required by the work plan.

\subsection{17. BA0-BA4 reopen review}\label{17-ba0-ba4-reopen-review}

\subsubsection{BA0}\label{ba0}

No authority/non-goal responsibility fails. Documentation remains
project authority.

\textbf{Disposition:} NOT TRIGGERED.

\subsubsection{BA1}\label{ba1}

Canonical keys/names do not require a third BAE identity family.

\textbf{Disposition:} NOT TRIGGERED.

\subsubsection{BA2}\label{ba2}

No new operator or role is materially forced by T3. The existing
extension rule remains adequate.

\textbf{Disposition:} NOT TRIGGERED.

\subsubsection{BA3}\label{ba3}

Cross-baseline rename and historical interpretation are handled by
existing continuity/provenance responsibilities plus immutable registry
revision references where applicable.

\textbf{Disposition:} NOT TRIGGERED.

\subsubsection{BA4}\label{ba4}

Shared rendering and method-owned interpretation remain separable. No
projection counterexample requires method taxonomy to become BA
vocabulary.

\textbf{Disposition:} NOT TRIGGERED.

\subsection{18. What T3 explicitly does not
claim}\label{18-what-t3-explicitly-does-not-claim}

T3 does not claim:

\begin{itemize}
\tightlist
\item
  that natural language has no synonyms;
\item
  that arbitrary legacy prose can be migrated automatically;
\item
  that semantic-kind vocabulary is universally complete;
\item
  that future projects can never force BA5 reopening;
\item
  that translations or fuzzy search are useless;
\item
  that an LLM cannot assist authoring;
\item
  that ThreatForge determines canonical vocabulary; or
\item
  that Base Analysis as a whole is closed.
\end{itemize}

It claims only that current DDTA evidence does not require normative
synonym machinery in the semantic core when authoring is controlled by
construction.

\subsection{19. Closure decision}\label{19-closure-decision}

The integrated T1/T2 contract survives the T3 adversarial review without
a material counterexample.

\begin{Shaded}
\begin{Highlighting}[]
\NormalTok{BA5{-}T3       CLOSED / PASS}
\NormalTok{BA5          CLOSED FOR CURRENT THESIS SCOPE}

\NormalTok{synonym support in normative core      REJECTED AS UNNECESSARY}
\NormalTok{optional lexical/NLP assistance        DEFERRED}
\NormalTok{BA0{-}BA4 reopen                         NOT TRIGGERED}
\NormalTok{BA6                                    NOT STARTED / NEXT PHASE}
\end{Highlighting}
\end{Shaded}

The final BA5 contract is recorded in the companion R1 contract artifact
included in this package.

\subsection{20. Next boundary}\label{20-next-boundary}

BA5 closure authorizes BA6 as the next phase, but this microstep does
\textbf{not} execute it.

BA6 must regress the complete BA0-BA5 design through the integrated
chain:

\begin{Shaded}
\begin{Highlighting}[]
\NormalTok{governed documentation}
\NormalTok{ {-}\textgreater{} canonical controlled authoring}
\NormalTok{ {-}\textgreater{} accepted BA}
\NormalTok{ {-}\textgreater{} BA3 provenance/change continuity}
\NormalTok{ {-}\textgreater{} multiple BA4 projections}
\NormalTok{ {-}\textgreater{} governed change}
\NormalTok{ {-}\textgreater{} rebuild / re{-}analysis}
\end{Highlighting}
\end{Shaded}

Only BA6 may close the complete Base Analysis milestone for the current
thesis scope.

\end{document}
